\clearpage
\section{Tra cứu nhanh}
\begingroup\small
\begin{multicols}{2}
\subsection{Dấu hiệu nhận dạng dạng bài}
Các dấu hiệu sau gợi ý hướng thử; luôn đối chiếu ràng buộc và
chứng minh điều kiện áp dụng.
\begin{itemize}
\setlength{\itemsep}{2pt}
\item \textbf{Tối ưu một giá trị số}: nếu mệnh đề đạt ít nhất $x$
      (hoặc không quá $x$) đơn điệu, thử tìm kiếm nhị phân trên đáp án.
      Ghi rõ hai biên và chiều đơn điệu.
\item \textbf{Nhiều truy vấn, dữ liệu cố định}: cân nhắc tiền xử lý
      bằng prefix sum, sparse table, LCA hoặc xử lý ngoại tuyến.
      Nếu có cập nhật, xác định phép gộp rồi xét Fenwick/segment tree.
\item \textbf{Phụ thuộc trên DAG}: topo rồi DP theo thứ tự topo.
      Nếu có chu trình, cần trạng thái hoặc lập luận khác.
\item \textbf{Mỗi đỉnh có đúng một cạnh đi ra}: đồ thị hàm gồm các
      chu trình và cây đổ vào chúng; tách chu trình khi đếm hoặc nhảy xa.
\item \textbf{Chọn/bỏ từng phần tử, $n$ nhỏ}: xét bitmask DP hoặc
      meet-in-the-middle. $2^{20}$ xấp xỉ $10^6$, nhưng còn phải tính
      số chuyển trạng thái và số bộ test.
\item \textbf{Ghép cặp hợp lệ}: kiểm tra cấu trúc hai phía.
      Nếu có chi phí, xét bài toán gán; nếu có sức chứa, xét luồng.
\item \textbf{Điều kiện kéo theo giữa các lựa chọn}: thử mô hình min-cut;
      cạnh sức chứa vô hạn cấm các cách chia không hợp lệ.
\item \textbf{Quá trình ngẫu nhiên}: DP xác suất/kỳ vọng trên trạng thái;
      nếu có vòng lặp, lập hệ phương trình và xét khả năng không dừng.
\item \textbf{Trò chơi hai người}: xác định thắng/thua từ các nước đi.
      Dùng Grundy khi trò chơi là tổng các phần độc lập phù hợp.
\item \textbf{Tọa độ nguyên, hỏi hướng/diện tích}: ưu tiên tích có hướng
      số nguyên; giao điểm và đường tròn cần xét sai số, suy biến.
\end{itemize}

\subsection{Gợi ý tư duy ad hoc và constructive}
\paragraph{Tìm cấu trúc.}
Thử dữ liệu nhỏ và cực trị; thay số cụ thể bằng thứ tự, chẵn lẻ,
phần dư hoặc quan hệ bằng nhau. Sắp xếp rồi xét đổi chỗ hai phần tử
kề nhau: nếu đổi chỗ không làm đáp án xấu hơn, chứng minh bằng
\emph{exchange argument} để suy ra dạng nghiệm chuẩn.
Với thao tác lặp, tìm bất biến hoặc thế năng nguyên bị chặn
và thay đổi chặt để chứng minh không lặp vô hạn.

\paragraph{Dựng và chứng minh.}
Trước hết tìm điều kiện cần bằng bất biến, giới hạn tổng, chẵn lẻ,
bậc hoặc đếm. Khi dựng từng bước, ghi rõ điều kiện còn giữ sau bước đó.
Chứng minh luôn có bước tiếp theo, quá trình kết thúc, và cấu hình cuối
thỏa \emph{mọi} điều kiện. Nếu lập luận chỉ đúng từ một kích thước nào đó,
xử lý riêng các trường hợp nhỏ. Với tham lam theo thứ tự từ điển,
phải chứng minh lựa chọn hiện tại vẫn cho phép hoàn thành phần còn lại.

\paragraph{Tìm phản ví dụ.}
Thử $n=0,1,2$ nếu hợp lệ; tất cả bằng nhau; tăng/giảm nghiêm ngặt;
giá trị âm, $0$, dương; các giá trị biên; nhiều nghiệm tối ưu.
Với tham lam, tìm hai lựa chọn đầu tiên dẫn đến kết quả khác nhau.
Vét cạn $n$ nhỏ giúp bác bỏ giả thuyết, nhưng không thay thế chứng minh.
\subsection{Công thức thường quên}
\begin{itemize}
\setlength{\itemsep}{2pt}
\item $\sum_{i=1}^{n}i=n(n+1)/2$;
      $\sum_{i=1}^{n}i^2=n(n+1)(2n+1)/6$;
      $\binom n2=n(n-1)/2$.
\item $1+r+\cdots+r^{n-1}=(r^n-1)/(r-1)$ khi $r\ne1$;
      với $r=1$, tổng bằng $n$. Khi tính modulo, chỉ chia cho $r-1$
      nếu nghịch đảo modulo tồn tại.
\item Với $P[0]=0$, $P[i+1]=P[i]+a[i]$, tổng nửa mở $[l,r)$
      là $P[r]-P[l]$. Độ dài giao hai đoạn nửa mở bằng
      $\max(0,\min(r_1,r_2)-\max(l_1,l_2))$.
\item Với $a\ge0$, $b>0$:
      $\lceil a/b\rceil=\lfloor a/b\rfloor+[a\bmod b\ne0]$.
      Cách viết $(a+b-1)/b$ có thể tràn số nguyên.
\item Với $a,b\ne0$:
      $\operatorname{lcm}(a,b)=|a/\gcd(a,b)\cdot b|$;
      chia trước khi nhân nhưng vẫn kiểm tra tràn.
\item Với $m>0$, phần dư chuẩn hóa của $a$ là
      $((a\bmod m)+m)\bmod m$ về mặt toán học.
      Trong C++, dùng \texttt{r=a\%m; if(r<0) r+=m;}
      để tránh cộng thêm $m$ khi không cần.
      Chia modulo $m$ chỉ hợp lệ khi mẫu nguyên tố cùng nhau với $m$.
\item Với $x>0$, $x\mathbin{\&}(-x)$ lấy bit $1$ thấp nhất;
      $x\mathbin{\&}(x-1)$ xóa bit $1$ thấp nhất.
\item $2S=|\sum_i\operatorname{cross}(P_i,P_{i+1})|$
      cho đa giác đơn. Khi chỉ so sánh khoảng cách, dùng bình phương
      khoảng cách để tránh \texttt{sqrt}.
\item $\Pr(A\cup B)=\Pr(A)+\Pr(B)-\Pr(A\cap B)$.
      Với $X$ nguyên không âm,
      $\mathbb E[X]=\sum_{k\ge1}\Pr(X\ge k)$,
      kể cả khi tổng có thể vô hạn.
\end{itemize}

\subsection{Checklist lỗi và trường hợp biên}
\paragraph{Trước khi cài.}
Xác định đoạn đóng hay nửa mở, chỉ số từ $0$ hay $1$, trạng thái vô
nghiệm, đồ thị có thể rời rạc/có cạnh song song/khuyên hay không.
Ước lượng thời gian và bộ nhớ ở ràng buộc lớn nhất, kể cả nhiều test.
\paragraph{Trước khi nộp.}
\begin{itemize}
\setlength{\itemsep}{2pt}
\item Kiểm tra $n=0$ hoặc $1$ nếu hợp lệ, tập rỗng, mọi phần tử bằng nhau,
      không có cạnh, đồ thị không liên thông, đáp án bằng $0$.
\item Kiểm tra tràn ở tích, tổng đường đi, $\mathrm{INF}+w$,
      tích ma trận, diện tích và $\operatorname{lcm}$;
      ép kiểu \texttt{1LL} hoặc \texttt{\_\_int128} \emph{trước} khi nhân.
\item Kiểm tra chia cho $0$, modulo âm, nghịch đảo không tồn tại;
      số thực gần $0$, đường song song/trùng, điểm trùng/thẳng hàng.
\item Xóa trạng thái giữa các test: mảng đánh dấu, danh sách kề,
      queue, biến toàn cục, nén tọa độ và đáp án.
\item Kiểm tra \texttt{i+1}, \texttt{i-1},
      \texttt{lower\_bound == end}, và phép trừ trên \texttt{size\_t}.
\item Nếu đệ quy sâu đến $O(n)$, xem giới hạn stack. Thử một test
      tự tính tay và vài test nhỏ ngẫu nhiên đối chiếu với vét cạn;
      khi sai, giữ seed và rút nhỏ phản ví dụ.
\end{itemize}
\subsection{Cú pháp Python cho số nguyên lớn}
Chỉ dùng nếu Python được phép và thời gian đủ. \texttt{int} không bị
giới hạn $64$ bit, nhưng phép toán trên số rất dài vẫn tốn thời gian.
\begin{itemize}
\setlength{\itemsep}{2pt}
\item \texttt{a = int(input())}; với nhiều dòng, dùng
      \texttt{sys.stdin.buffer.readline}. Xuất bằng \texttt{print(a)}.
\item \texttt{a // b} lấy sàn kể cả khi $a<0$;
      \texttt{a \% b} bằng $0$ hoặc cùng dấu với $b$ khi $b\ne0$.
      \texttt{divmod(a,b)} trả về thương và dư.
\item \texttt{pow(a,k,m)} tính $a^k\bmod m$ mà không tạo số $a^k$
      đầy đủ ($k\ge0$, $m>0$). \texttt{pow(a,-1,m)} tìm nghịch đảo
      nếu $\gcd(a,m)=1$; nếu không sẽ phát sinh \texttt{ValueError}.
\item \texttt{math.comb(n,k)}, \texttt{math.gcd(a,b)} và
      \texttt{math.isqrt(n)} cho tổ hợp, UCLN và
      $\lfloor\sqrt n\rfloor$ chính xác; \texttt{isqrt} cần $n\ge0$.
      \texttt{comb}, \texttt{isqrt} và số mũ âm của \texttt{pow}
      với modulo cần Python 3.8 trở lên.
\item Trên Python 3.11 trở lên, chuyển đổi số thập phân cực dài có thể
      vướng giới hạn chữ số. Khi đề thật sự cần, kiểm tra
      \texttt{sys.get\_int\_max\_str\_digits()} và có thể gọi
      \texttt{sys.set\_int\_max\_str\_digits(0)}.
      Kiểm tra phiên bản Python của máy chấm.
\end{itemize}

\subsection{trick lỏ}
pragma lấy trên codeforces ok
\lstinputlisting[language=C++,basicstyle=\ttfamily\scriptsize]{code/quick_reference/pragmas.txt}
\end{multicols}
\endgroup
